\documentclass[10pt,a4paper]{amsart}
\usepackage[margin=19mm]{geometry}
\usepackage{amsmath,amssymb,mathtools}
\usepackage[scr=rsfs,cal=euler]{mathalfa}
\usepackage{xfrac}
\usepackage[unicode,hidelinks,bookmarks=false]{hyperref}
\newcommand{\ie}{i.e.\ }
\newcommand{\cf}{cf.\ }
\newcommand{\resp}{respectively\ }
\newcommand{\Subsection}{\subsection}
\providecommand{\nobreakdash}{\nobreak\hbox{-}}
\setlength{\emergencystretch}{2em}
\multlinegap0pt
\usepackage{amsthm}
\usepackage{mathtools}
\usepackage{amsmath,amssymb}
\usepackage{enumitem}
\usepackage{graphicx}
\def\veca{{\mbox{\boldmath \(a\)}}}
\def\vecb{{\mbox{\boldmath \(b\)}}}
\def\vecm{{\mbox{\boldmath \(m\)}}}
\newtheorem{theorem}{Theorem}[section]
\newtheorem{corollary}[theorem]{Corollary}
\newtheorem{lemma}[theorem]{Lemma}
\newtheorem{fact}[theorem]{Fact}
\newtheorem{example}{Example}
\newtheorem{proposition}[theorem]{Proposition}
\title[Optimization and complexity of inertia-type bounds on the independence number]{Optimization and Complexity of Inertia-Type Bounds on the Independence Number: the equivalence of the two mixed-integer formulations (Theorem 3.2)}
\author[A. Abiad, S. van Hoesel, V. Michaux]{Aida Abiad, Stan van Hoesel, Valentin Michaux}
\date{Source: arXiv:2605.00524v2 (CC BY 4.0)}
\begin{document}
\maketitle
\vspace{1em}
\noindent\textit{Source and licence.} Optimization and Complexity of Inertia-Type Bounds on the Independence Number. Version arXiv:2605.00524v2 (CC BY 4.0), distributed under the Creative Commons Attribution 4.0 International licence, http://creativecommons.org/licenses/by/4.0/, as recorded in the arXiv OAI licence metadata for this exact version and shown on the abstract page https://arxiv.org/abs/2605.00524v2. This item is an editorial re-presentation of the paper's formulation-equivalence theorem together with the paper's own proof of it. A full rights notice is delivered at licenses/arxiv-2605.00524-CC-BY-4.0.txt.
\noindent\textit{Editorial scope.} Counted unit: the paper's equivalence theorem for its two mixed-integer formulations of the $k$-independence number, printed as Theorem 3.2 (source0.tex lines 437--448), delivered together with the paper's complete original proof of it (source0.tex lines 450--490) as the single primary proof body. No mathematics is written, repaired or re-derived: statement and proof are the paper's own text line for line, in the paper's own notation ($z_{1,u}$, $z_2$, $\textsc{MILP}\,1(u)$, $\textsc{MILP}\,2$, $V(G)$ and the decision vectors of the two formulations). The two mixed-integer formulations themselves are defined earlier in the paper (their labels lie outside the retained spans); the counted proof needs no cross-reference to them, because the paper states the transformation argument in full in terms of the decision vectors that the delivered proof introduces. Every cross-reference inside the retained spans resolves inside the delivered document, and no citation command occurs in any retained span, so the delivered document carries no bibliography and no reference or citation command of the delivered text was rewritten or adapted. Printed numbers are the paper's own: the wrapper restores the paper's shared section-relative theorem counter before the counted statement, so it prints as Theorem 3.2. Omitted from this package and not redistributed: the paper's preamble and front matter as such (of which only the title and the author names are reproduced in the wrapper); the introduction, the two mixed-integer formulations with their polyhedral and normalization discussion, the remaining theorems and lemmas of the paper, the complexity analysis and the conclusion; and the paper's whole bibliography with its bib data file. Printed slips kept literal, among them the paper's own wording and its own choice of symbols. Layout and class: the paper's own class is the standard LaTeX article class at 10pt, which is not a publisher class; the delivered wrapper is the lane's pinned amsart editorial wrapper at 10pt on A4, to which this package adds the paper's own theorem-like declarations with their shared section-relative counter and the shorthand macros its retained text uses, all copied verbatim from the paper's source. No publisher class file, no publisher style file and no upstream script is redistributed. Assets: no retained span of this package contains a figure environment or a graphics-inclusion command, no image or artwork file exists in or is redistributed with this package, and the manifest and the delivery evidence both carry an empty asset-dependency list. A full rights notice is delivered at licenses/arxiv-2605.00524-CC-BY-4.0.txt.
\setcounter{section}{3}
\setcounter{theorem}{1}
\begin{theorem}
\label{thm:milp-equivalence}
Let \(z_{1,u}\) be the optimal value of \(\textsc{MILP}\,1(u)\) for
\(u\in V(G)\), and let \(z_2\) be the optimal value of
\(\textsc{MILP}\,2\). Then
\[
    z_2=\min_{u\in V(G)}z_{1,u}.
\]
Moreover, every feasible solution of \(\textsc{MILP}\,2\) can be transformed
into a feasible solution of \(\textsc{MILP}\,1(u)\) for some \(u\in V(G)\),
without changing its binary variables or objective value.
\end{theorem}

\begin{proof}
Fix \(u\in V(G)\) and let \((\veca,\vecb)\) be feasible for
\(\textsc{MILP}\,1(u)\). Its diagonal constraints give
\((p(A))_{uu}=0\) and \((p(A))_{vv}\geq0\) for every
\(v\neq u\). Hence, all diagonal constraints of \(\textsc{MILP}\,2\) are
satisfied. Since the spectral and integrality constraints are identical in
the two formulations, \((\veca,\vecb)\) is feasible for
\(\textsc{MILP}\,2\). The objective is also unchanged, as both formulations
minimize \(\vecm^\top\vecb\). Therefore,
\(z_2\leq z_{1,u}\) for every \(u\in V(G)\), and thus
\(z_2\leq\min_{u\in V(G)}z_{1,u}\).

Conversely, let \((\veca,\vecb)\) be feasible for
\(\textsc{MILP}\,2\), let \(p(x)=\sum_{i=0}^{k}a_i x^i\), and define
\(\gamma=\min_{v\in V(G)}(p(A))_{vv}\). By the diagonal constraints,
\(\gamma\geq0\). Choose \(\widehat u\in V(G)\) such that
\((p(A))_{\widehat u\widehat u}=\gamma\), and define
\(\widetilde p(x)=p(x)-\gamma\). Equivalently, only the constant coefficient
is changed, from \(a_0\) to \(\widetilde a_0=a_0-\gamma\).

Since \(\widetilde p(A)=p(A)-\gamma I\), we have
\((\widetilde p(A))_{\widehat u\widehat u}=0\) and \((\widetilde p(A))_{vv}\geq0\) for every \(v\in V(G)\setminus\{\widehat u\}\). Thus, the diagonal
constraints of \(\textsc{MILP}\,1(\widehat u)\) are satisfied. Furthermore,
for every \(j=0,\ldots,d\),
\(\widetilde p(\theta_j)=p(\theta_j)-\gamma\leq p(\theta_j)\). Hence,
\(p(\theta_j)-b_j+\varepsilon\leq0\) implies
\(\widetilde p(\theta_j)-b_j+\varepsilon\leq0\), so the spectral constraints
remain valid.

The transformation modifies only the continuous coefficient \(a_0\), while
the binary vector \(\vecb\) is kept unchanged. Since
\((\veca,\vecb)\) is feasible for \(\textsc{MILP}\,2\), it satisfies
\(\vecb\in\{0,1\}^{d+1}\); therefore,
\((\widetilde{\veca},\vecb)\) satisfies the same integrality constraint.
Moreover, both formulations minimize \(\vecm^\top\vecb\), so keeping
\(\vecb\) unchanged preserves the objective value. Consequently,
\((\widetilde{\veca},\vecb)\) is feasible for
\(\textsc{MILP}\,1(\widehat u)\) with the same objective value, which yields
\(\min_{u\in V(G)}z_{1,u}\leq z_2\). Combining both inequalities proves the
claim.
\end{proof}

\end{document}
